\section{Introduction}

Model predictive control (MPC) repeatedly queries a dynamics model while
optimizing an action sequence. Its performance therefore depends on two
coupled quantities: the state at the forecast origin and the transition used
to propagate that state under candidate controls \cite{hewing2020_learning_mpc}.
For a learned transition, improving one quantity does not guarantee improving
the other. A model may forecast accurately from clean states yet initialize
poorly from incomplete measurements, or it may correct the current state while
accumulating transition error during rollout.

World models provide trainable predictors for planning and continuous control
\cite{hansen2023_td_mpc2_scalable_robust_world,hafner2025_mastering_diverse_control_tasks_through}.
Physics-guided models add a nominal transition or physical constraint, which
can reduce data demand and expose departures from the assumed mechanism
\cite{karniadakis2021_physics_informed_machine_learning,yu2024_learning_dynamical_systems_from_data}.
Neither construction by itself resolves how recent noisy or missing
observations should be reconciled with nominal dynamics before planning.

This paper studies that interface through the Complementary-Filtered Residual
World Model (\method). During an eight-step observed context, nominal dynamics
predict the next state and a mask-gated complementary update combines the
prediction with available measurements. The final corrected state initializes
a probabilistic nominal-physics-plus-neural-residual transition. Future states
are then rolled out freely under a candidate control sequence. The design
separates forecast-origin correction from transition learning and can be
queried directly by an optimizer.

The evaluation connects predictive accuracy to realized regulation. The
open-loop study contains five nonlinear controlled systems, three data regimes,
eight independently generated data/training seeds, 11 model variants, and
\ConfirmatoryRuns{} trained models. A frozen additional study places four
prediction interfaces inside the same cross-entropy MPC controller and
evaluates \ClosedLoopEpisodes{} paired control episodes. An external experiment
uses measured NanoDrone states \cite{busetto2026_nanodrone} with a deliberately
reduced quadrotor prior to test transfer beyond the synthetic equations.

The contributions are:
\begin{enumerate}
  \item a causal estimator--predictor architecture that combines a mask-gated
  complementary context update, nominal dynamics, and a bounded neural
  residual for free multistep rollout;
  \item an error decomposition and constrained mean-square-optimal scalar gain
  that identify when a fixed complementary update should help;
  \item a seed-paired evaluation that links open-loop forecast error to
  receding-horizon regulation without retraining or controller-specific model
  selection; and
  \item component, outage, calibration, and measured-flight analyses that show
  where the context filter helps and where the residual and predictive
  intervals do not.
\end{enumerate}

The results do not reduce to a single accuracy claim. The aggregate open-loop
contrast favors \method, whereas the clean-regime gate fails; the fixed
gain-0.5 filter followed by nominal rollout outperforms the learned residual on
aggregate point error; and the measured-flight study separates position error
from full-state error and interval coverage. These distinctions determine how
the model should be used in an automation pipeline.

